% sec-A069881.tex -- k = 2 section of the unified kp+1 paper.
% Standalone fragment, \input by A083833-UNIFIE-paper.tex.
% Environments and notation inherited from sec-carry.tex; the
% guards below are no-ops once defined there.
\providecommand{\IsPal}{\mathrm{IsPal}}
\makeatletter
\@ifundefined{theorem}{\newtheorem{theorem}{Theorem}}{}
\@ifundefined{corollary}{\newtheorem{corollary}[theorem]{Corollary}}{}
\makeatother

\section{Multiplier $k = 2$: A069881}
\label{sec:k2}

Entry A069881~\cite{oeis} (A.~Murthy, April 2002) lists
$1, 2, 3, 4, 5, 55, 151, 161, 171, 181, 191, 252, \ldots$
Its only comment (J.~McCann, 2002) records that every
repdigit $55\ldots5$ is a term, and the classification
of even-length members is the conjecture of
\emph{Tangente} problem~8154~\cite{tangente} quoted in
Section~\ref{sec:intro}. Instantiating the carry chain
of Section~\ref{sec:carry} at $k = 2$ gives
$c_j = \bigl\lfloor\bigl(2\,(p \bmod 10^j) + 1\bigr)
/10^j\bigr\rfloor$, hence $c_0 = 1$ and $c_j \le 1$.

\begin{theorem}[Tangente 8154, full form]
\label{thm:k2-even}
Let $m \ge 1$ be even and $\IsPal(p,m)$. Then
$2p+1$ is a palindrome (of some length $m'$) if and
only if every digit of $p$ equals $5$, i.e.\ $p$ is the
repdigit $5\,R_m$.
\end{theorem}

\begin{theorem}[member dichotomy]
\label{thm:k2-dicho}
If $\IsPal(p,m)$ and $\IsPal(2p+1,m')$ with $m \ge 1$,
then either every digit of $p$ equals $5$ (the overflow
branch, where $2p+1 = 11\ldots1$ has $m+1$ digits), or
$m$ is odd, $a_0 \in \{1,2,3,4\}$, and $a_j \le 4$ at
even positions while $a_j \ge 5$ at odd positions.
\end{theorem}

\begin{corollary}[image parity]
\label{cor:k2-img}
Under the same hypotheses, every digit of $q = 2p+1$ at
an even position is odd.
\end{corollary}

\begin{proof}
Same-length branch ($m' = m$):
Lemma~\ref{lem:same} at $k = 2$ forces $m$ odd,
$a_0 \in \{1,2,3,4\}$ and the displayed classes.
Overflow branch ($m' = m+1$): the boundary
$(2a_0 + 1) \bmod 10 = c_m = 1$ forces $a_0 = 5$ and
$c_1 = 1$; at each subsequent position the mirrored
carry is forced to $1$ by parity ($2a'+1$ is odd),
while the joint constraints give $a' \in \{0,5\}$ and
$a' \ge 5$: the chain is rigid in a single step, so
$a_j = 5$ throughout, $p = 55\ldots5$ and
$q = 11\ldots1$. The converse direction of
Theorem~\ref{thm:k2-even} is the kernel-verified
identity $2 \cdot 5R_m + 1 = R_{m+1}$ packaged as a
literal iff (Table~\ref{tab:lean}); the forward
direction follows because even $m$ excludes the
same-length branch.
\end{proof}

The statements are consistent with all $45$ published
terms: the only even-length ones are the repdigits $55$
and $5555$, and every other listed term has odd length,
units digit in $\{1,2,3,4\}$ and the alternating classes
of Theorem~\ref{thm:k2-dicho}.
